\gls{CRC} is a type of tumor that affects the large intestine and the rectum, which are responsible for processing food waste and absorbing water and nutrients. After evaluating the symptoms, the diagnosis is confirmed by exams such as a fecal occult blood test, \gls{CL}, and a biopsy. For a successful treatment, early diagnosis of the disease is essential. However, the absence of symptoms in the early stages, associated with a lack of access to screening exams, contributes to late diagnoses.

The goal of this work is to use \glspl{DNN} for the early detection of \gls{CRC} using biopsy images. The proposed methodology involves the comparison of different architectures, from classic ones like AlexNet to those utilizing residual connections like ResNet. Furthermore, it includes the early adoption of quantum convolutional layers in \gls{CRC} classification.

The results obtained demonstrate how the use of \gls{ML} can be useful for the early detection of \gls{CRC}, contributing to more accurate diagnoses and treatments, thereby increasing the chances of a cure and improving the quality of life for patients.

